\selectlanguage{american}%
Volume computation of a given set is a fundamental problem in theoretical
computer science, with diverse applications to other problems (see
Exercise~\ref{ex: linExt-volume}). For a convex body $K$ in $\R^{n}$,
$R$ such that $B_{n}\subseteq K\subseteq RB_{n}$ and $\varepsilon>0$,
the problem is to output a number $A$ s.t. 
\[
(1-\varepsilon)\vol(K)\le A\le(1+\varepsilon)\vol(K).
\]

More precisely, we have the following basic problem. 
\begin{problem}[Volume Estimation]
\label{problem:volume} Given a membership oracle for a convex body
$K$ in $\R^{n}$, $R$ such that $B_{n}\subseteq K\subseteq RB_{n}$
and $\varepsilon>0$, output a number $A$ s.t. 
\[
(1-\varepsilon)\vol(K)\le A\le(1+\varepsilon)\vol(K).
\]

The complexity of an algorithm for volume estimation is measured by
the number of oracle queries and the number of arithmetic operations. 
\end{problem}


\section{Hardness of Volume Estimation}
\begin{thm}
\cite{BF87,BF88} Suppose there is a deterministic algorithm that
takes a convex body $K\subset\R^{n}$ given by a membership oracle
and outputs $A(K),B(K)$ such that $A(K)\le\vol(K)\le B(K)$ and makes
at most $n^{a}$ calls to the membership oracle for $K$ for some
$a>0$. Then there is some convex body $K$ for which

\[
\frac{B(K)}{A(K)}\ge\left(\frac{cn}{a\log n}\right)^{\frac{n}{2}}
\]
where $c>0$ is an absolute constant.
\end{thm}

We will prove the following simpler bound.
\begin{thm}
\cite{E86} For any deterministic algorithm that makes $m$ queries
to a membership oracle for a convex body $K\subseteq B^{n}$ and returns
$A(K),B(K)$ such that $A(K)\le\vol(K)\le B(K)$, there is some convex
body $K$ for which 
\[
\frac{B(K)}{A(K)}\ge\frac{2^{n}}{m}.
\]
\end{thm}

In other words, even getting a simple exponential volume approximation
requires an exponential number of queries. 
\begin{proof}
For all queries $x\in B^{n},$the oracle answers YES and for queries
outside the unit ball it answers NO. After $m$ queries $x_{1},\ldots,x_{m}$,
we have that 
\[
\conv\left\{ x_{1},\ldots,x_{m}\right\} \subseteq K\subseteq B^{n}.
\]
Now by Lemma , the volume of the convex hull of \emph{any $m$ }points
in the unit ball is at most a $m/2^{n}$ fraction of the unit ball's
volume. This proves the theorem.
\end{proof}
\begin{lem}
Let $x_{1},\ldots,x_{m}\in B^{n}$. Then 
\[
\vol\left(\conv\left\{ x_{1},\ldots,x_{m}\right\} \right)\le\frac{m}{2^{n}}\vol(B^{n}).
\]
\end{lem}

\begin{proof}
For each $x_{i}$, define $B_{i}=B(x_{i},\|x_{i}\|/2)$, a ball with
diameter $\|x_{i}\|\le1/2$. Then 
\[
\vol\left(\bigcup_{i=1}^{m}B_{i}\right)\le m\cdot\vol(B(0,1/2))=\frac{m}{2^{n}}\vol(B^{n}).
\]

We argue that $\conv{x_{1},\ldots,x_{m}}\subset\bigcup B_{i}$. To
see this, let $y$ be a point not contained in $B_{i}$. Then, the
angle at $y$ induced by the vectors $y$ and $y-x_{i}$ must be acute,
i.e., 
\[
\langle y,x_{i}-y\rangle<0
\]
and this holds for all $i$. In other words, the halfspace $\langle y,x-y\rangle<0$
separates $y$ from all the $x_{i}$. Therefore it separates it from
their convex hull. 
\end{proof}
Computing the volume exactly is hard even for explicit simple polyhedra
as the following theorem asserts.
\begin{thm}
\cite{DF88} Let $P={x:x\in[0,1]^{n},a^{\top}x\le b}$ where $a$
is an integer vector. Computing the volume of $P$ (specified by $a,b$)
is $\#$P-hard. 
\end{thm}


\section{Cutting Plane method for Volume Computation: Reduction to Center-of-Gravity}

The cutting plane method provides a simple algorithm for computing
the volume of a bounded set or, more generally, the integral of a
function. The algorithm below relies crucially on the ability to compute
the center of gravity (or centroid) of the original set/function restricted
by halfspaces. For simplicity we describe it here for the case when
the input object is a convex body. This reduction from volume computation
to centroid computation is deterministic, and hence establishes the
hardness of the latter problem.

\begin{algorithm2e}[H]

\caption{$\mathtt{CuttingPlaneVolume}$}\label{alg:cutting-plane-volume}

\SetAlgoLined

\textbf{Input:} A body $K\subset\R^{n}$, $r>0$, and an oracle for
computing the centroid.

Fix an ordering on the axes: $e_{1},e_{2},\ldots e_{n}$.

$K^{(0)}=K,$ $z^{(0)}=\mbox{centroid}(K^{(0)}),k=0,V_{0}=1$.

\For{$k=0,\cdots$}{

Let $e_{i}$ be an axis vector so that the width of $K^{(k)}$ along
$e_{i}$ is greater than $r$.

\lIf{There is no such $e_{i}$}{\textbf{Break}.}

Let $a$ be $e_{i}$ or $-e_{i}$ with sign chosen so that $H^{(k)}\defeq\left\{ x\,:\,a^{\top}x\le a^{\top}z^{(k)}\right\} $
contains the origin.

$K^{(k+1)}\leftarrow K^{(k)}\cap H^{(k)}$.

$z^{(k+1)}\leftarrow\mbox{centroid}(K^{(k+1)})$.

$\widehat{z}\leftarrow\mbox{centroid}(K^{(k)}\setminus K^{(k+1)})$.

$V_{k+1}=V_{k}\cdot\frac{\norm{\widehat{z}-z^{(k+1)}}}{\norm{\widehat{z}-z^{(k)}}}$.

}

Return $V_{\text{last}}\cdot\prod_{i=1}^{n}w_{i}(K^{(k+1)})$ where
$w_{i}(\cdot)$ is the width along $e_{i}$.

\end{algorithm2e}

The idea behind the algorithm is simple: when we cut a set with a
hyperplane through its centroid, the line joining the centroids of
the two sides passes through the original centroid (see Figure~\ref{fig:cutting-volume});
moreover, the ratio of the two segments is exactly the ratio of the
volumes of the two halfspaces (i.e., $V_{k}=\vol(K^{(0)})/\vol(K^{(k)})$).
\begin{lem}
\label{lem:centroid}For any measurable bounded set S in $\R^{n}$
and any halfspace $H$ with bounding hyperplane containing the centroid
of $S$, we have 
\[
\mbox{centroid}(S)=\frac{\vol(S\cap H)}{\vol(S)}\,\mbox{centroid}(S\cap H)+\frac{\vol(S\cap\overline{H})}{\vol(S)}\,\mbox{centroid}(S\cap\overline{H}).
\]
\end{lem}

The following lemma has a proof similar to that of the Grunbaum theorem.
\begin{lem}
\label{lem:width}Let $K\subseteq\R^{n}$ be a convex body with centroid
at the origin. Suppose that for some unit vector $\theta$, the support
of $K$ along $\theta$ is $[a,b]$. Then, 
\[
|a|\ge\frac{b}{n}.
\]
\end{lem}

\begin{proof}
The goal of this proof is to show that a cone attains the maximum
ratio of $b/|a|$ (and it equals $n$ in that case). As with the proof
of Theorem \ref{thm:Grunbaum}, we may assume $K$ is symmetric by
replacing each cross section $\{x\in K\colon\theta^{\top}x=t\}$ with
the $(n-1)$-dimensional ball of the same volume with radius $r(t)$,
i.e., $r^{n-1}(t)C_{n-1}$ is the $(n-1)$-dimensional volume of $K\cap\{y\in\R^{n}:y^{\top}\theta=t\}$,
with the center of gravity at $t=0$. We will exploit the fact that
$r(t)$ is concave in $t$, and replace $r(t)$ by a line going through
$(b,0)$ and $(0,r(0))$ so that we get a cone with base at $t=a$.
We will show that the center of gravity of the cone, say $t_{0}$,
is to the left of $0$. Then, 
\[
\frac{b}{|a|}\le\frac{b-t_{0}}{t_{0}-a}=\frac{n}{n+1}\frac{n+1}{1}=n
\]
which can be rearranged to give the result. The equation for the cone
is
\[
y(t)=-\frac{r(0)}{b}(x-b).
\]
Notice that $y(t)\le r(t)$ for all $t\in[0,b]$ since it is a chord
between two points of a concave function and therefore always below
the concave function. To see why $r(t)\le y(t)$ for $t\in[-a,0]$,
consider $g(t)=r(t)-y(t)$. Then $g'(t)=r'(t)-y'(t)$. Notice that
$g'(0)=r'(0)-y'(0)\ge0$, and since $y'(t)$ is a constant and $r'(t)$
is decreasing, we must have $r'(t)-y'(t)\ge0$ for all $t\le0$. Since
$g(0)=0$ by construction of $y(t)$, we get that $g(t)\le0$ for
all $t\le0$. Then, we have
\begin{align*}
t_{0} & =\frac{\int_{a}^{b}ty^{n-1}(t)dt}{\int_{a}^{b}y^{n-1}(t)dt}\\
 & =\frac{1}{\int_{a}^{b}y^{n-1}(t)dt}\left(\int_{a}^{0}ty^{n-1}(t)dt+\int_{0}^{b}ty^{n-1}(t)dt\right)\\
 & \le\frac{1}{\int_{a}^{b}y^{n-1}(t)dt}\left(\int_{a}^{0}tr^{n-1}(t)dt+\int_{0}^{b}tr^{n-1}(t)dt\right)=0.
\end{align*}
So $t_{0}\le0$.
\end{proof}
Using the above property, we can show that the algorithm reaches a
cuboid in a small number of iterations.
\begin{thm}
Let $K$ be a convex body in $\R^{n}$ such that $[-r,r]^{n}\subseteq K\subseteq[-R,R]^{n}$.
Algorithm $\mathtt{CuttingPlaneVolume}$ correctly computes the volume
of $K$ using $O(n\log\frac{nR}{r})$ centroid computations.
\end{thm}

\begin{proof}
By Theorem \ref{thm:Grunbaum}, at each iteration, the volume of the
remaining set $K^{(k)}$ decreases by at least a factor of $(1-\frac{1}{e})$.
When the directional width along any axis is less than $r$ (namely
$\max_{x\in K}\left|e_{i}^{\top}x\right|\leq r$), the algorithm stops
cutting along that axis. So, just before the last cut along any axis,
the width in that direction is at least $r$. Then, we use the center
of gravity to cut. By Lemma \ref{lem:width}, the directional width
along every axis of the surviving set is at least $r/(n+1)$. Since
the set always contains the origin, and the original set contains
a cube of side length $r$, when the algorithm stops, it must be an
axis-parallel cuboid with each side of length in the range $[r/(n+1),r]$.
So the final volume is at least $(r/(n+1))^{n}.$ The initial volume
is at most $R^{n}$. Therefore the number of iterations is at most
\[
\log_{(1-\frac{1}{e})}\left(\frac{(2R)^{n}}{(r/(n+1))^{n}}\right)=O\left(n\log(nR/r)\right).
\]
In each iteration, by Lemma \ref{lem:centroid}, the algorithm maintains
the ratio of the volume of the original $K$ to the current $K^{(k)}$.

\selectlanguage{english}%
\begin{figure}
\centering{}\includegraphics[width=0.5\columnwidth]{Figures/Section_9\lyxdot 2}\caption{\foreignlanguage{american}{Algorithm (\ref{alg:cutting-plane-volume}) in its second iteration. }}
\end{figure}
\end{proof}
The above algorithm shows that computing the volume is polytime reducible
to computing the centroid. Since volume is known to be \#P-hard for
explicit polytopes, this means that centroid computation is also \#P-hard
for polytopes \cite{rademacher2007approximating}. In later chapters
we will see randomized polytime algorithms for sampling and hence
for approximating centroid and volume. 
\begin{xca}
\cite{brightwell1991counting} \label{ex: linExt-volume} Given a
partial order $P$ on an $n$-element set, it is of interest to count
the number of linear extensions, i.e., total orders on the set that
are consistent with $P$. We can define an associated polyhedron,
\[
Q=\left\{ x\in[0,1]^{n}:\,x_{i}<x_{j}\mbox{ if }(i,j)\in P\right\} .
\]
Show that the number of linear extensions of $P$ is exactly $\vol(Q)n!$. 
\end{xca}


\section{Proof of the Brunn-Minkowski Inequality\label{sec:Proof-of-BM}}

Here we prove Theorem \ref{thm:Brunn-Minkowski}. Let $A,B\subset\R^{n}$
be measurable subsets. Our goal is to show that 
\[
\vol(\lambda A+(1-\lambda)B)^{1/n}\ge\lambda\vol(A)^{1/n}+(1-\lambda)\vol(B)^{1/n}.
\]
By setting $A'=\lambda A$ and $B'=(1-\lambda)B$, this is equivalent
to showing that for all measurable $A,B$:
\begin{equation}
\vol(A+B)^{1/n}\ge\vol(A)^{1/n}+\vol(B)^{1/n}.\label{eq:BM-ineq}
\end{equation}

We first consider the case when $A,B$ are cuboids, i.e., 
\[
A=\prod_{i=1}^{n}[0,a_{i}]\ \mbox{and \ }B=\prod_{i=1}^{n}[0,b_{i}]
\]
for positive real numbers $a_{1},\ldots,a_{m},b_{1},\ldots,b_{m}$.
Then, we have 
\[
\vol(A)=\prod_{i=1}^{n}a_{i},\,\vol(B)=\prod_{i=1}^{n}b_{i}\:\vol(A+B)=\prod_{i=1}^{n}(a_{i}+b_{i}).
\]
Now we need to verify that 
\[
\left(\prod_{i=1}^{n}(a_{i}+b_{i})\right)^{\frac{1}{n}}\ge\left(\prod_{i=1}^{n}a_{i}\right)^{\frac{1}{n}}+\left(\prod_{i=1}^{n}b_{i}\right)^{\frac{1}{n}}
\]
which is equivalent to 
\[
\left(\prod_{i=1}^{n}\frac{a_{i}}{a_{i}+b_{i}}\right)^{\frac{1}{n}}+\left(\prod_{i=1}^{n}\frac{b_{i}}{a_{i}+b_{i}}\right)^{\frac{1}{n}}\le1.
\]
Using the AM-GM inequality, for each of the terms on the LHS, we have
\[
\left(\prod_{i=1}^{n}\frac{a_{i}}{a_{i}+b_{i}}\right)^{\frac{1}{n}}+\left(\prod_{i=1}^{n}\frac{b_{i}}{a_{i}+b_{i}}\right)^{\frac{1}{n}}\le\frac{1}{n}\sum_{i=1}^{n}\frac{a_{i}}{a_{i}+b_{i}}+\frac{1}{n}\sum_{i=1}^{n}\frac{b_{i}}{a_{i}+b_{i}}=1.
\]

We will now prove the inequality to any finite union of cuboids. That
suffices to complete the proof because any measurable subset is the
limit of a union of disjoint cubes as the cubes become smaller and
smaller. 

Let $A$ be a finite collection of disjoint axis-parallel cuboids
and $B$ be another such collection. We will prove the inequality
by induction on the total number of cuboids in $A\cup B$. First note
that there is at least one complete cuboid of $A$ separated from
another complete cuboid of $A$ by an axis-aligned halfspace, i.e.,
there is a coordinate $i$ so that after translating $A$ along this
coordinate the sets ${x:x_{i}\le0}$and ${x:x_{i}>0}$ both contain
a complete cuboid of $A$. This follows from the fact that disjoint
axis-parallel cuboids can be separated by an axis-parallel halfspace.
Now define the following sets:
\begin{align*}
A^{+} & =A\cap\left\{ x:x_{i}>0\right\} \ \ A^{-}=A\setminus A^{+}\\
B^{+} & =B\cap\left\{ x:x_{i}>0\right\} \ \ B^{-}=B\setminus B^{+}.
\end{align*}
Translate $B$ so that we have 
\begin{equation}
\frac{\vol(A^{+})}{\vol(A)}=\frac{\vol(B^{+})}{\vol(B^{-})}.\label{eq:AB_ratio}
\end{equation}
(Translation only translates the Minkowski sum.) Since $A^{+}\cup B^{+}$
and $A^{-}\cup B^{-}$ are both smaller collections of cuboids compared
to $A\cup B$, we can apply the B-M inequality (\ref{thm:Brunn-Minkowski})
to them inductively. We already proved the base case when each of
$A,B$ is a single cuboid. Using induction in the second step and
the identity (\ref{eq:AB_ratio}) in the fourth step, we have (note
that the sets $A^{+}+B^{+}$and $A^{-}+B^{-}$ are disjoint and contained
in $A+B$):
\begin{align*}
\vol(A+B) & \ge\vol(A^{+}+B^{+})+\vol(A^{-}+B^{-})\\
 & \ge\left(\vol(A^{+})^{1/n}+\vol(B^{+})^{1/n}\right)^{n}+\left(\vol(A^{-})^{1/n}+\vol(B^{-})^{1/n}\right)^{n}\\
 & =\vol(A^{+})\left(1+\frac{\vol(B^{+})^{1/n}}{\vol(A^{+})^{1/n}}\right)^{n}+\vol(A^{-})\left(1+\frac{\vol(B^{-})^{1/n}}{\vol(A^{-})^{1/n}}\right)^{n}\\
 & =\vol(A^{+})\left(1+\frac{\vol(B)^{1/n}}{\vol(A)^{1/n}}\right)^{n}+\vol(A^{-})\left(1+\frac{\vol(B)^{1/n}}{\vol(A)^{1/n}}\right)^{n}\\
 & =\vol(A)\left(1+\frac{\vol(B)^{1/n}}{\vol(A)^{1/n}}\right)^{n}\\
 & =\left(\vol(A)^{1/n}+\vol(B)^{1/n}\right)^{n}
\end{align*}
which completes the proof.\selectlanguage{english}%

